\section{Annotation Protocol and Prompt Templates}
\label{app:annotation}
\label{app:prompts}

\paragraph{Policy label annotation.}
For datasets without gold policy labels (WhoQA and FaithDial), we assigned action labels from $\mathcal{A} = \{\textsc{Answer}, \textsc{Disclaimer}, \textsc{Disclose}, \textsc{Abstain}\}$ using the following procedure. Two annotators (the paper authors) independently labeled a shared sample of 200 examples using the rubric below. Cohen's $\kappa$ was 0.71 on this shared sample. The remaining development examples were split evenly and labeled by a single annotator each, with cross-check sampling at 10\% to detect drift.

\noindent \textbf{Annotation rubric:}
\begin{itemize}
\item \textsc{Answer}: At least two retrieved passages directly support a consistent answer, with no contradicting passage.
\item \textsc{Disclaimer}: Only one passage supports the answer, or the supporting passage contains hedged language (``may'', ``reportedly'', ``some sources suggest'').
\item \textsc{Disclose}: At least two passages make directly contradictory factual claims about the question's target entity.
\item \textsc{Abstain}: No retrieved passage is topically relevant to the question, or the question asks for a specific fact that none of the passages address.
\end{itemize}

\paragraph{RAG prompt template.}
The following prompt is used for the \textsc{Answer} action:
\begin{quote}
\small
\texttt{You are a factual question-answering assistant. Use ONLY the following passages to answer the question. If the passages do not contain enough information to answer, say ``I cannot answer based on the provided passages.''}\\
\texttt{Passages: \{passages\}}\\
\texttt{Question: \{question\}}\\
\texttt{Answer:}
\end{quote}

\paragraph{GPT-4 Verifier prompt.}
The GPT-4 Verifier uses the following structured prompt:
\begin{quote}
\small
\texttt{You will classify how a RAG system should respond given the following question and retrieved passages. Choose one action from: ANSWER, DISCLAIMER, DISCLOSE, ABSTAIN. Return a JSON object with fields: action, rationale (one sentence), confidence (0.0--1.0).}\\
\texttt{Question: \{question\}}\\
\texttt{Passages: \{passages\}}
\end{quote}
